\section{Methodology}
\label{sec:method}

\subsection{Overview}

The key insight driving our design is that adversarial vulnerability, viewed as a scalar, discards
the multivariate structure that carries architecture-family signal. We therefore represent each
model as a $\Delta^*$-vector --- one normalized vulnerability score per attack category --- and
test whether architecture family explains a significant portion of variance in this multivariate
space.

\subsection{The $\Delta^*$-Vector Framework}

\paragraph{Normalized Vulnerability ($\Delta^*$).}
For a model $m$ evaluated on attack category $c$:
\begin{equation}
\Delta^*(m, c) = \frac{\text{Acc}_\text{clean}(m, c) - \text{Acc}_\text{adv}(m, c)}{\text{Acc}_\text{clean}(m, c)}
\end{equation}
The normalization by $\text{Acc}_\text{clean}$ removes the clean-accuracy confound: $\Delta^*$
measures \emph{proportional} degradation from each model's own baseline, isolating the adversarial
vulnerability effect from task difficulty and model capacity. For each model, $\Delta^*$ scores
across $K$ reliable attack categories form a vector $\mathbf{x}(m) \in \mathbb{R}^K$. The data
matrix $\mathbf{X} \in \mathbb{R}^{N \times K}$ ($N=9$, $K=6$) is the input to statistical
analysis.

\paragraph{Reliability Filtering.}
For each attack category $c$, we compute the Spearman--Brown corrected split-half reliability
$r_\text{SB}$. Categories with $r_\text{SB} < 0.7$ or fewer than 50 examples are excluded.
This filter makes effect-size estimates conservative: only stable attack categories enter the
analysis.

\subsection{Model Selection and Scale Matching}

Nine transformer models spanning three families, all in the 110--350M parameter range:

\begin{table}[h]
\centering
\caption{Model pool by architecture family.}
\begin{tabular}{ll}
\toprule
\textbf{Family} & \textbf{Models} \\
\midrule
Encoder-only & bert-base-uncased, roberta-base, \\
             & google/electra-base-discriminator, albert-base-v2 \\
Decoder-only & gpt2, facebook/opt-125m, facebook/opt-350m \\
Encoder-decoder & t5-base, facebook/bart-base \\
\bottomrule
\end{tabular}
\end{table}

Scale matching controls for the capacity confound; we do not include 1B+ models in this PoC.

\subsection{Fine-Tuning Protocol}

All models fine-tuned on GLUE tasks (SST-2, MNLI, QQP, QNLI, RTE) using HuggingFace Trainer
with AdamW optimizer, family-specific learning rates (encoder-only: $2 \times 10^{-5}$;
decoder-only: $5 \times 10^{-5}$; encoder-decoder: $1 \times 10^{-4}$), and 10\% linear warmup.
Seed=42 throughout.

\textbf{Note on enc\_dec coverage:} T5-base and BART-base were fine-tuned on SST-2 only (not
MNLI) in this experiment, producing degenerate ANLI-R3 results for the enc\_dec family
(see Section~\ref{sec:discussion}).

\subsection{Statistical Methods}

\paragraph{Permutation MANOVA.}
We test whether architecture family membership explains $\Delta^*$-vector variance using
permutation MANOVA on $\mathbf{X}$ with family labels $\mathbf{y}$. Test statistic: Pillai's
trace; effect size: $\eta^2 = \text{SS}_\text{between}/\text{SS}_\text{total}$. In the balanced
three-group case, Pillai's trace and $\eta^2$ are monotonically related (both bounded [0,1]) but
not identical --- we report $\eta^2$ throughout for interpretability, with Pillai's trace as the
basis for the permutation $p$-value. $P$-value: empirical fraction of 1,000 permuted $\eta^2$
exceeding observed $\eta^2$. Permutation approach is valid at $N=9$ without distributional
assumptions. We additionally compute per-category $\eta^2$ via univariate ANOVA for decomposition.

\paragraph{LOMO Classification.}
Leave-one-model-out nearest-neighbor (cosine distance, $k=1$) classification. Reports accuracy
and $3 \times 3$ confusion matrix. \emph{Validity caveat:} at $N=3$ models per family, LOMO is
geometrically near-degenerate; results are exploratory only.

\paragraph{Mixed-Effects Controls.}
Sensitivity check with formula
$\Delta^*(m,c) \sim \text{arch\_family} \times \text{attack\_type} + \text{objective} +
\text{tokenizer} + \text{clean\_acc} + (1|\text{model\_id})$.
Note: at $N=9$, this model is severely underdetermined (collinearity/separation); results are
reported with appropriate caveats.

\subsection{Implementation}

Seven modules: \texttt{data\_loader.py}, \texttt{fine\_tuner.py}, \texttt{evaluator.py},
\texttt{delta\_star.py}, \texttt{statistical\_analysis.py}, \texttt{visualizer.py},
\texttt{run\_experiment.py}. Pipeline checkpoints results as JSON after each stage for crash
recovery.
